The electric power grid has undergone a fundamental transformation over the last two
decades. The integration of digital communications, programmable logic controllers, and
IP-connected SCADA and Human-Machine Interface (HMI) systems has enabled
unprecedented operational efficiency but has simultaneously expanded the cyber attack
surface. High-profile incidents—most notably the Stuxnet worm \cite{langner2011stuxnet}
and the Ukrainian power grid attacks of 2015 and 2016—demonstrate that adversaries are
willing and able to traverse the IT/OT boundary and compromise energy infrastructure.

Regulatory frameworks respond to this threat. The North American Electric Reliability
Corporation Critical Infrastructure Protection (NERC CIP) standards mandate event logging,
monitoring, and incident response for Bulk Electric Systems. IEC~62351 specifies security
for power systems communications, including requirements for audit trail generation.
Together they establish Windows Event Logs—produced natively by every Windows host in a
smart grid operations centre—as a mandated, machine-readable record of system behaviour.

Despite this regulatory mandate, Windows Event Logs have received comparatively little
attention as a detection surface for ICS-specific attacks. The research community has focused
on network-level anomalies (e.g., Modbus/TCP deviations \cite{goldenberg2013modbus}) or
dedicated ICS protocols, leaving host-based event logs underexplored. Meanwhile, general
system-log anomaly detection has advanced significantly through deep learning approaches
such as DeepLog \cite{du2017deeplog} and LogAnomaly \cite{zhao2019loganomaly}, and
through log-parsing methods like Drain \cite{he2017drain}. These techniques have not been
applied in the ICS/SCADA context.

This paper makes the following contributions:
\begin{itemize}
  \item A synthetic Windows Event Log generator that models four ICS-relevant attack
        patterns within a realistic smart grid topology.
  \item A session-level feature engineering pipeline tailored to the temporal structure of
        Windows Event Logs.
  \item A baseline evaluation of Random Forest classification under stratified
        cross-validation, establishing measurable success criteria.
  \item An open-source framework enabling reproducible ICS security research without
        access to operational data.
\end{itemize}

Section~\ref{sec:lit} reviews related work. Section~\ref{sec:proposed} states the central
hypothesis and research objectives. Section~\ref{sec:method} describes the methodology.
Section~\ref{sec:results} reports preliminary results, and Section~\ref{sec:conc} concludes.
